\section{Application: plant disease diagnosis and forecasting}\label{sec:plant}

Plant disease diagnosis and forecasting use both arrows of the diagram every day. A
diagnostician reasons from a causal hypothesis to what must be observed, which is deduction, and
from symptoms, signs and assay results to the hypotheses they leave open, which is induction. A
forecast reasons from weather and field conditions to the risk of disease and is revised season by
season. Unlike \cref{sec:quantum}, this application needs no mathematics beyond
\cref{sec:setting,sec:deduction,sec:induction} and one new definition, the resolution of an
assay (\cref{def:resolves}). Diagnosis and disease management decisions have long been studied
quantitatively in plant disease epidemiology~\cite{hughes1999,madden2007}; the aim here is to
place them in the framework, not to replace those methods.

\begin{table}[ht]
\centering\small
\begin{tabular}{@{}l>{\raggedright\arraybackslash}p{0.62\textwidth}@{}}
\toprule
Symbol & Plant pathology \\
\midrule
$\W$ (worlds) & candidate states of a plant or field: causal agent (if any), host, conditions \\
$\varphi$, $t$ & a diagnosis (``the cause is $X$'') or a disease-risk model \\
$\exists$ & a symptom, a sign, or the result of an assay \\
$s$ (study) & a diagnostic work-up, or the record of past seasons \\
$\models$ (hence) & the hypothesis predicts the observation \\
$\circlearrowright$ (conclude) & the diagnosis remains consistent with the record \\
$i$ (iteration) & a further assay, or a further season \\
$h$ (hypothesis index) & the choice of assay (\cref{def:resolves}) \\
$\up,\down$ & a positive or negative result; disease forecast or not \\
\bottomrule
\end{tabular}
\caption{Plant-pathology reading of the symbols of \cref{tab:symbols}.}
\label{tab:plant}
\end{table}

\subsection{Elimination and refinement}

The deductive results of \cref{sec:deduction} and the refinement results of
\cref{sec:induction} apply without change. If a diagnosis $H$ entails an observation $o$ and $o$
fails, $H$ is eliminated (\cref{prop:tollens}); this is why one clean negative result can end a
candidate diagnosis. Each further result can only shrink the set of diagnoses that remain
supported (\cref{thm:mono}), so a diagnosis once eliminated returns only if an earlier result is
withdrawn.

\subsection{Assay resolution}

\begin{definition}[Assay and resolution]\label{def:resolves}
An \emph{assay} is a map $a:\W\to R$ to a set of results $R$, and the observation ``$a$ gave
$r$'' is the thought $[a=r]$, which holds in $w$ when $a(w)=r$. The assay $a$ \emph{resolves} a
diagnosis $H$ if worlds with the same result agree on $H$:
\[
  a(w)=a(w')\ \Longrightarrow\ \bigl(H(w)\iff H(w')\bigr)\qquad\text{for all }w,w'\in\W.
\]
\lean{AssayResult, Resolves}
\end{definition}

\begin{theorem}[Assay resolution]\label{thm:resolves}
An assay $a$ resolves $H$ if and only if, for every world $w$, the result $[a=a(w)]$ entails $H$
or entails $\neg H$.
\lean{resolves\_iff}
\end{theorem}

\begin{proof}
Suppose $a$ resolves $H$ and fix $w$. If $H(w)$, every $v$ with $a(v)=a(w)$ satisfies $H$, so
$[a=a(w)]\models H$; otherwise the same argument gives $[a=a(w)]\models\neg H$. Conversely, let
$a(w)=a(w')$. Then $w$ and $w'$ both satisfy $[a=a(w)]$; if this result entails $H$ both satisfy
$H$, and if it entails $\neg H$ neither does.
\end{proof}

\begin{corollary}[Underdetermination by an assay]\label{cor:undetermined}
If $a(w)=a(w')$, $H(w)$ and $\neg H(w')$, then the result $[a=a(w)]$ entails neither $H$ nor
$\neg H$. If $a$ resolves $H$, so does the joint assay $w\mapsto(a(w),b(w))$ for any further
assay $b$.
\lean{undetermined, Resolves.pair\_left}
\end{corollary}

\begin{proof}
The world $w'$ satisfies the result but not $H$, and $w$ satisfies it and $H$. For the joint
assay, equal pairs have equal first components.
\end{proof}

Within the model, repeating an assay adds an observation that holds in exactly the worlds where the
first one holds, so it removes no world. The limit of \cref{cor:undetermined} therefore lies in what
the assay can tell apart, not in how often it is run. This gives the hypothesis index $h$ of the
foundation a concrete reading: $h$ selects the assay, and the assay bounds what any study can
reveal.

\begin{example}[Formae speciales]\label{ex:fusarium}
Pathogenic strains of \emph{Fusarium oxysporum} are classified into formae speciales by host
specificity~\cite{jackson2024}: f.\,sp.\ \emph{cubense} causes Fusarium wilt (Panama disease) of
banana, and f.\,sp.\ \emph{lycopersici} causes Fusarium wilt of tomato. Take a two-world model in
which culture morphology gives the same result for both, and a host test records the host on
which the isolate causes wilt. By \cref{cor:undetermined}, morphology does not resolve ``the
isolate is f.\,sp.\ \emph{cubense}'', however many cultures are examined; by
\cref{thm:resolves}, the host test does. The same pattern covers races distinguished only on
differential hosts and cryptic species distinguished only by molecular markers.
\lean{morphology\_not\_resolves, hostTest\_resolves}
\end{example}

\Cref{cor:undetermined} is the classical form of \cref{prop:qhume}(4): measurements in one basis
do not resolve the phase of a spin, just as morphology does not resolve the forma specialis. What
is specifically quantum is that no single basis is informationally complete for all states,
whereas classically the assay that reports the world itself resolves every hypothesis.

\subsection{Koch's postulates as hypothetico-deductive confirmation}

Koch's postulates establish that a pathogen causes a disease by isolating it from diseased
plants, inoculating it into healthy hosts, reproducing the symptoms and re-isolating the same
organism~\cite{agrios2005}. The inoculation step has the form of \cref{thm:hd}.

\begin{proposition}[Koch's postulates]\label{prop:koch}
Let $H$ be ``$X$ causes the disease'' and $o$ be ``inoculating healthy hosts with $X$ reproduces
the symptoms''. Suppose $H$ is possible, $H\models o$, and some world falsifies $o$. Then, under
every prior on finitely many worlds, $o$ confirms $H$; and if $o$ fails, $H$ is eliminated.
\lean{confirms\_of\_entails, modus\_tollens}
\end{proposition}

\begin{proof}
The first claim is \cref{thm:hd} with the study $(o)$, the second \cref{prop:tollens}.
\end{proof}

The condition that $o$ could fail excludes tests that succeed whatever the cause; mock-inoculated
controls serve this purpose in practice, by showing that the procedure itself does not produce the
symptoms. The entailment $H\models o$ holds only under conditions favourable to disease: a
susceptible host, enough inoculum and a suitable environment, the disease triangle of plant
pathology~\cite{agrios2005}. A failed inoculation therefore eliminates $H$ only together with
these auxiliary assumptions.

\subsection{Diagnostic tests and forecasts}

Diagnostic tests and forecasts are usually judged by counts. Of $D$ diseased and $N$ healthy
units (plants, samples or fields), suppose $tp$ diseased and $fp$ healthy units test positive, or
are forecast positive.

\begin{proposition}[When a positive result raises the risk]\label{prop:plantconfirm}
The share of diseased units among the positives exceeds the prevalence exactly when the
sensitivity exceeds the false-positive rate:
\[
  \frac{tp}{tp+fp}>\frac{D}{D+N}\quad\iff\quad\frac{tp}{D}>\frac{fp}{N}.
\]
In cross-multiplied form, $D(tp+fp)<tp\,(D+N)\iff D\,fp<tp\,N$ for all natural numbers, with no
positivity assumption.
\lean{positive\_raises\_risk\_iff}
\end{proposition}

\begin{proof}
Expand both products: $D\,tp+D\,fp<tp\,D+tp\,N$, and cancel $D\,tp$.
\end{proof}

\Cref{prop:plantconfirm} is \cref{def:confirm} in frequency form: a positive result confirms
``diseased'' exactly when the test discriminates. The prevalence affects how far the risk rises,
not whether it rises, as the prior does in \cref{cor:bswan}. Whether the higher risk justifies
action is a separate decision, in which the risk is compared with a threshold set by the costs of
treatment and of loss~\cite{hughes1999,madden2007}.

A forecast is an assay applied before disease appears, to weather and field variables. For
Sclerotinia stem rot of soybean, weather-based models predict whether apothecia of
\emph{Sclerotinia sclerotiorum}, the source of ascospore inoculum, are present during flowering,
with 67--70\% accuracy on independent data sets~\cite{willbur2018}. Three results of the framework
apply to such forecasts.
\begin{enumerate}
  \item A risk model fitted to past seasons is supported by them but never entailed: the next
        season is an unobserved individual, so \cref{thm:hume} applies with seasons in place of
        birds.
  \item \Cref{prop:plantconfirm} applies with ``forecast positive'' as the test.
  \item Forecasting rules are defeasible, as the next proposition shows.
\end{enumerate}

\begin{proposition}[A defeasible forecasting rule]\label{prop:sclerotinia}
Consider a soybean field at flowering under weather favourable to disease, with worlds
$(\mathrm{noApothecia},\mathrm{infected})\in\{0,1\}^2$ ranked so that apothecia with infection
is most normal, no apothecia without infection next, and the other two combinations least normal.
Then the field is normally infected, $()\mathrel{|\!\sim}\mathrm{infected}$, but
$(\mathrm{noApothecia})\not\mathrel{|\!\sim}\mathrm{infected}$; indeed
$(\mathrm{noApothecia})\mathrel{|\!\sim}\neg\mathrm{infected}$.
\lean{sclerotinia\_nonmonotone, noApothecia\_normally\_not\_infected}
\end{proposition}

The ranking is a stylized model: infection without local apothecia, for instance from ascospores
blown in from neighbouring fields, is ranked abnormal rather than impossible. Like the Tweety
example of \cref{prop:tweety}, the forecast is closed under the rules of \cref{thm:klm} and is
withdrawn when new information arrives.

The decision thresholds, the accuracy figures and the biological premises of the models above are
empirical; the Lean formalization covers the logical results only.
